\section{Introduction}

Tensor decompositions provide compact models for multiway data in scientific
data analysis, signal processing, and masked multiway learning
\citep{Kolda2009Tensor,Sidiropoulos2017TSP,Liu2013TensorCompletion,
LiuMoitra2020TensorCompletion}. Tucker decomposition is a
standard choice because it separates a small core from mode-wise factors. It
uses a linear map from the Tucker latent tensor to the measured tensor: the
reconstructed Tucker value is fitted directly to the measurement. This
assumption is natural for many data sets,
but it can be mismatched to measurement or data-generation pipelines with
response curves, saturation, calibration effects, or material-dependent
nonlinearities
\citep{Debevec1997HDR,Grossberg2004CameraResponse,
BioucasDias2012HyperspectralUnmixing,Heylen2014NonlinearUnmixing,
Luo2024LRTFR}. In such settings, the latent signal may be well organized by a
low-rank Tucker tensor while the measured tensor appears to have larger
multilinear rank after a nonlinear response.

A direct remedy is to increase Tucker rank, but this changes the multilinear
budget and expands the core and factors throughout the model. More flexible
nonlinear tensor decoders based on kernels, Gaussian processes, or neural
networks can also be useful, but they address a broader problem and often make
the relation to a Tucker backbone less explicit
\citep{xu2011infinite,zhe2016distributed,liu2017deepcp,
saragadam2024deeptensor,ahn2024neural}. We ask a narrower question: can a
Tucker model keep its chosen backbone ranks fixed while learning a response
that explains coherent nonlinear residuals left by matched-rank Tucker?

We instantiate this idea with a polynomial link. Let
\[
\cS = \tucker{\cX}{\bA^{(1)}, \ldots, \bA^{(N)}}
\]
be a Tucker latent tensor. Instead of fitting the measured tensor directly by
\(\cS\), we model
\[
\widehat{\cY} = f_{\bbeta}(\cS) = \sum_{p=0}^{P} \beta_p \cS^{\odot p},
\]
where \(\odot p\) denotes the Hadamard power. We call the model Nonlinear Tucker Decomposition
with Polynomial Links (NTD-PL). It contains Tucker as the degree-one case,
keeps the Tucker backbone and multilinear ranks fixed, and adds only \(P+1\)
scalar link coefficients. The higher powers \(\cS^{\odot p}\) induce
high-order interactions among the same core and factors, so the added
expressivity is parameter-tied to one learned latent tensor.

This viewpoint is central to the paper. NTD-PL is best understood as a Tucker
enhancement and shared-backbone rank lift, not as a generic nonlinear decoder. A
degree-\(p\) power of the latent tensor behaves like an ordinary Tucker tensor
whose factors are polynomial feature lifts of the original factors, but those
lifted features remain tied to one learned backbone. The new degrees therefore
reshape the low-rank geometry through the polynomial link, rather than
adding an untied high-rank tensor block. This is also why the link is not a
post-hoc calibration layer.

Our claims are intentionally scoped. Increasing Tucker rank remains the right
comparison when the residual is ordinary multilinear structure, and we do not
claim to dominate neural, kernel, or other nonlinear tensor factorization
methods. The intended regime is more specific: the residual left by
matched-rank Tucker should be coherent with a nonlinear response.

This scoped regime is also supported empirically. In controlled tensors, the gain grows with the injected nonlinear residual energy and with the degree needed to approximate the response. On real hyperspectral tensors, NTD-PL improves matched-rank Tucker in both full reconstruction and held-out completion. More importantly, a one-step Link Refresh diagnostic predicts when NTD-PL will help: across HSI scenes and non-HSI tensor units, larger diagnostic reduction is associated with larger NTD-PL gains, while low-yield units show little or no benefit. These results suggest that nonlinear responses are identifiable structure in several real tensor families.

\paragraph{Contributions.}
We make the following contributions.
\begin{itemize}[leftmargin=*,itemsep=2pt,topsep=2pt]
  \item \textbf{Model.} We propose NTD-PL, a Tucker enhancement with a
  polynomial link on a Tucker latent tensor. It keeps the Tucker
  backbone and multilinear ranks fixed while adding only scalar coefficients
  for a parameter-tied response.
  \item \textbf{Theory.} We characterize the induced model class through
  Veronese rank expansion, separation from matched-rank Tucker, and
  approximation results for nonlinear responses.
  \item \textbf{Optimization.} We derive a masked objective for reconstruction
  and completion, and optimize it with Link Refresh, Linked-Backbone Descent,
  and degree continuation.
  \item \textbf{Experiments.} We evaluate controlled nonlinear residuals,
  hyperspectral reconstruction and completion, rank-lift and Link Refresh
  diagnostics, and external cross-domain boundary cases.
\end{itemize}
